\documentclass{article}
\usepackage{amsmath, amssymb, amsthm}
\usepackage{hyperref}
\usepackage{geometry}
\geometry{margin=1in}

\title{Erd\H{o}s Problem \#956}
\date{Last updated: 2025-08-31}
\author{Source: erdosproblems.com}

\begin{document}
\maketitle

\noindent\textbf{Status:} OPEN \quad \textbf{No prize}

\noindent\textbf{Tags:} geometry, distances, convex

\medskip
\noindent\textbf{Problem Statement:}

If $C,D\subseteq \mathbb{R}^2$ then the distance between $C$ and $D$ is defined by\[\delta(C,D)=\inf_{\substack{c\in C\\ d\in D}}\| c-d\|.\]Let $h(n)$ be the maximal number of unit distances between disjoint convex translates. That is, the maximal $m$ such that there is a compact convex set $C\subset \mathbb{R}^2$ and a set $X$ of size $n$ such that all $(C+x)_{x\in X}$ are disjoint and there are $m$ pairs $x_1,x_2\in X$ such that\[\delta(C+x_1,C+x_2)=1.\]Determine $h(n)$ - in particular, prove that there exists a constant $c>0$ such that $h(n)>n^{1+c}$ for all large $n$.

\bigskip
\noindent\textbf{Remarks:}

A problem of Erd\H{o}s and Pach \cite{ErPa90}, who proved that $h(n) \ll n^{4/3}$. They also consider the related function where we consider $n$ disjoint convex sets (not necessarily translates), for which they give an upper bound of $\ll n^{7/5}$. 

It is trivial that $h(n)\geq f(n)$, where $f(n)$ is the maximal number of unit distances determined by $n$ points in $\mathbb{R}^2$ (see [90]).

\bigskip
\noindent\textbf{References:}

[ErPa90] Erd\H{o}s, P. and Pach, J., Variations on the theme of repeated distances. Combinatorica (1990), 261--269.

\bigskip
\noindent\small{Source: \url{https://www.erdosproblems.com/956}}
\end{document}
